\section{Data Fusion}
\label{sec:data_fusion}

Data fusion merges all records that describe the same real-world entity into a single consolidated record. If attribute values from different records conflict with each other, a conflict resolution method is applied to resolve the conflict by selecting or generating a single value~\cite{bleiholder2009data, naumann2006dataFusionThreeSteps}. 

\textbf{Conflict Resolution.}\label{sec:fusion_approach} 
The PyDI data integration framework offers a range of different conflict resolution methods including voting, numeric aggregations (e.g., average, median), string-based methods (e.g., longest string, shortest string), temporal (e.g., most recent), and source-aware methods (e.g., source priority based on estimated source accuracy). To perform data fusion, one of these conflict resolution methods is selected for each attribute in the target schema.

\textbf{Human Pipeline.} Graduate students first created validation sets of 10--23 entities per use case. To select entities, they browsed entity clusters containing conflicting attribute values and chose entities whose correct values could be reliably verified through external sources. For each selected entity and each conflicting attribute, they performed web searches to determine the ground truth value, consulting authoritative sources such as official company websites, music databases, or game registries. They then assigned a conflict resolution method to each attribute by choosing from PyDI's library of resolvers. For each attribute, the students considered the data type and domain semantics to select an appropriate method, for instance median for numeric scores and revenue. They iteratively tested different conflict resolution methods by comparing fused values to the corresponding values in the validation set and chose the configuration resulting in the highest value overlap.

\textbf{LLM Pipeline.} The LLM pipeline also generates fusion validation sets that are used to choose conflict resolution methods. The pipeline generated fusion validation sets in two steps: First, GPT-5.2 is prompted to select well-known entities from a random sample of 100 entity groups, i.e., entities for which it is likely easy to gather ground truth (e.g., The Beatles in Music, Apple in Companies). GPT-5.2 selects approximately 30 groups per use case. Second, the model is prompted to generate the value that it perceives as correct for each conflicting attribute. We test two variants: LLM-only validation, where the model generates values from its parametric knowledge, and RAG validation, where the LLM verifies its answers against web search results. The pipeline queries the OpenAI web search tool for creating the RAG validation set.

Using these validation sets, the pipeline evaluates multiple candidate configurations generated by combining heuristic and LLM-based rule selection strategies with iterative refinement, and selects the configuration with the highest validation accuracy (see footnote~\ref{fn:prompts} for all prompts).

\textbf{Results.}\label{sec:fusion_results} Table~\ref{tab:fusion_results} compares fusion test accuracy on held-out test sets of 15--23 entities. The first column shows the accuracy of the students' manually configured fusion strategies (Human Config). The remaining columns show the accuracy of the best pipeline-generated configuration selected using human-created, LLM-only, or RAG validation sets.

% Data fusion results comparing validation set variants

\begin{table}[t]
\caption{Data fusion test accuracy by validation set origin.}
\label{tab:fusion_results}
\centering
\setlength{\tabcolsep}{4pt}
\begin{tabular}{@{}lcccc@{}}
\toprule
 & \textbf{Human} & \textbf{Human} & \textbf{LLM} & \textbf{LLM} \\
\textbf{Use Case} & \textbf{Config} & \textbf{Val.\ Set} & \textbf{Val.\ Set} & \textbf{+ RAG} \\
\midrule
Games     & .832 & .866 & .866 & .866 \\
Companies & .803 & .861 & .709 & .744 \\
Music     & .766 & .721 & .708 & .708 \\
\midrule
\textit{Average} & .800 & .816 & .761 & .773 \\
\bottomrule
\end{tabular}
\end{table}


The pipeline with human validation sets outperforms the manual human configuration on average (.816 vs.\ .800), indicating that automated candidate generation finds better resolver assignments even when using the same validation data. For Games, all three automated variants select the same configuration (.866), surpassing the manual configuration (.832). For Companies, the human validation set leads to the highest accuracy overall (.861), while LLM-only achieves .709 and RAG achieves .744. This gap is largely attributable to temporal mismatches: time-sensitive attributes such as headquarters city and revenue are difficult for the LLM to validate, as the source data reflects a specific point in the past while the LLM returns current values. For Music, LLM-only and RAG validation both select the same configuration (.708), close to the human validation set (.721), while the manual configuration achieves .766. On average, LLM-only validation achieves 76.1\% compared to 81.6\% for human validation. RAG validation provides a modest improvement to 77.3\%, suggesting that GPT-5.2's parametric knowledge is largely sufficient for well-known entities.

\textbf{Discussion.}\label{sec:fusion_discussion}
The results show that LLM-generated validation sets can guide fusion configuration selection comparably to human-created ones for domains with predominantly static attributes (Music, Games). The accuracy gap in Companies is largely attributable to temporal mismatches rather than fundamental limitations of the approach. Providing the LLM with metadata about the temporal context of each source could address this.

The selection of well-known entities for validation introduces a bias, as fusion strategies optimized for famous entities may not generalize to less prominent ones. The RAG variant is therefore most relevant not for the use cases studied here, but for enterprise scenarios involving internal or proprietary data sources where entities may not appear in the LLM's training data. In such settings, grounding validation in web or intranet search could substantially improve validation set quality.

